% !TEX program = xelatex
% ============================================================
%  Arabic Critical Path Method (CPM) Template — Nibras Center
%  مشروع بحثي: مخطط المسار الحرج
%  Compile with: xelatex main.tex (twice)
% ============================================================

\documentclass[11pt,a4paper,landscape]{article}

\usepackage[a4paper,landscape,margin=2cm,top=2.5cm,bottom=2.5cm,headheight=15pt]{geometry}
\usepackage{amsmath}
\usepackage{array}
\usepackage{booktabs}
\usepackage{tabularx}
\usepackage{enumitem}
\usepackage{float}
\usepackage{titlesec}
\usepackage{fancyhdr}
\usepackage{setspace}
\usepackage[table]{xcolor}
\usepackage{tikz}
\usetikzlibrary{positioning,calc,shapes.geometric,arrows.meta}

\usepackage{bidi}
\usepackage[no-math]{fontspec}
\usepackage[quiet]{polyglossia}

\setmainfont[Scale=1]{Amiri-Regular.ttf}
\newfontfamily\arabicfont[Script=Arabic,Scale=0.95]{Amiri-Regular.ttf}
\newfontfamily\englishfont[Scale=0.95]{TeX Gyre Termes}

\setdefaultlanguage[calendar=gregorian,hijricorrection=1]{arabic}
\setotherlanguage[variant=british]{english}

\usepackage[breaklinks,hidelinks]{hyperref}
\hypersetup{pdftitle={مخطط المسار الحرج},pdfauthor={الباحث الرئيسي}}

\titleformat{\section}{\Large\bfseries}{\thesection}{0.7em}{}
\titlespacing*{\section}{0pt}{1.2ex plus 0.3ex}{0.8ex plus 0.2ex}

\pagestyle{fancy}
\fancyhf{}
\fancyhead[R]{\small\itshape مخطط المسار الحرج (CPM)}
\fancyhead[L]{\small اسم المشروع}
\fancyfoot[C]{\small\thepage}
\renewcommand{\headrulewidth}{0.3pt}

\newcolumntype{L}[1]{>{\raggedright\arraybackslash}p{#1}}
\newcolumntype{C}[1]{>{\centering\arraybackslash}p{#1}}

\definecolor{criticalColor}{HTML}{C0392B}
\definecolor{normalColor}{HTML}{2E6DA4}
\definecolor{nodeFill}{HTML}{F0F4F8}
\definecolor{criticalFill}{HTML}{FADBD8}

\setstretch{1.3}

\begin{document}

\begin{center}
{\LARGE\bfseries مخطط المسار الحرج (CPM)}\\[0.5em]
{\Large عنوان المشروع البحثي}\\[1em]
\end{center}

% ============================================================
% 1. TASK LIST WITH DURATIONS
% ============================================================
\section{قائمة المهام والمدد}
\label{sec:tasks}

\begin{table}[H]
\centering
\small
\renewcommand{\arraystretch}{1.3}
\begin{tabularx}{\textwidth}{C{1cm} L{5cm} C{2cm} C{3cm} C{2cm} C{2cm}}
\toprule
\textbf{الرمز} & \textbf{المهمة} & \textbf{المدة (أسبوع)} & \textbf{المهام السابقة} & \textbf{أولى بدء} & \textbf{أواخر بدء} \\
\midrule
A & مراجعة الأدبيات & 6 & --- & 0 & 0 \\
B & تصميم المنهجية & 4 & A & 6 & 6 \\
C & تجهيز المعدات & 3 & A & 6 & 7 \\
D & جمع البيانات & 10 & B, C & 10 & 10 \\
E & التحليل الإحصائي & 6 & D & 20 & 20 \\
F & النمذجة والمحاكاة & 6 & D & 20 & 21 \\
G & التحقق من النتائج & 4 & E, F & 26 & 27 \\
H & كتابة الورقة & 8 & G & 30 & 31 \\
I & تقديم للنشر & 4 & H & 38 & 39 \\
J & التقرير النهائي & 6 & G & 30 & 30 \\
\bottomrule
\end{tabularx}
\end{table}

% ============================================================
% 2. NETWORK DIAGRAM
% ============================================================
\section{مخطط الشبكة (Network Diagram)}
\label{sec:network}

\begin{figure}[H]
\centering
\begin{tikzpicture}[
  every node/.style={font=\small},
  normal/.style={fill=nodeFill,draw=normalColor,thick,rounded corners=3pt,inner sep=6pt,minimum width=2cm,align=center},
  critical/.style={fill=criticalFill,draw=criticalColor,thick,rounded corners=3pt,inner sep=6pt,minimum width=2cm,align=center},
  normalline/.style={draw=normalColor,thick,-{Stealth[length=5mm]}},
  criticalline/.style={draw=criticalColor,very thick,-{Stealth[length=5mm]}},
]

% Nodes — critical path in red
\node[critical] (A) at (0,0) {A\\مراجعة الأدبيات\\(6)};
\node[critical] (B) at (4,0) {B\\تصميم المنهجية\\(4)};
\node[normal]  (C) at (4,-2.5) {C\\تجهيز المعدات\\(3)};
\node[critical] (D) at (8,0) {D\\جمع البيانات\\(10)};
\node[critical] (E) at (12,0) {E\\التحليل الإحصائي\\(6)};
\node[normal]  (F) at (12,-2.5) {F\\النمذجة\\(6)};
\node[critical] (G) at (16,0) {G\\التحقق\\(4)};
\node[normal]  (H) at (20,1.5) {H\\كتابة الورقة\\(8)};
\node[normal]  (I) at (24,1.5) {I\\تقديم للنشر\\(4)};
\node[critical] (J) at (20,-1.5) {J\\التقرير النهائي\\(6)};

% Edges
\draw[criticalline] (A) -- (B);
\draw[normalline]   (A) -- (C);
\draw[normalline]   (C) -- (D);
\draw[criticalline] (B) -- (D);
\draw[criticalline] (D) -- (E);
\draw[normalline]   (D) -- (F);
\draw[normalline]   (F) -- (G);
\draw[criticalline] (E) -- (G);
\draw[normalline]   (G) -- (H);
\draw[normalline]   (H) -- (I);
\draw[criticalline] (G) -- (J);

% Critical path label
\node[font=\bfseries\color{criticalColor}] at (12,-4) {المسار الحرج: A → B → D → E → G → J (36 أسبوع)};

\end{tikzpicture}
\caption{مخطط الشبكة للمسار الحرج — المسار الحرج باللون الأحمر}
\end{figure}

% ============================================================
% 3. CRITICAL PATH ANALYSIS
% ============================================================
\section{تحليل المسار الحرج}
\label{sec:analysis}

\textbf{المسار الحرج:} \color{criticalColor}\textbf{A → B → D → E → G → J}\color{black}\\[0.5em]
\textbf{مدة المشروع الإجمالية:} \textbf{36 أسبوع} (9 أشهر)

\begin{table}[H]
\centering
\caption{تحليل المسار الحرج}
\renewcommand{\arraystretch}{1.3}
\begin{tabularx}{\textwidth}{C{1cm} L{4cm} C{2cm} C{2cm} C{2cm} X}
\toprule
\textbf{الرمز} & \textbf{المهمة} & \textbf{أولى بدء} & \textbf{أواخر بدء} & \textbf{السماحية} & \textbf{حرجة؟} \\
\midrule
A & مراجعة الأدبيات & 0 & 0 & 0 & \color{criticalColor}\textbf{نعم} \\
B & تصميم المنهجية & 6 & 6 & 0 & \color{criticalColor}\textbf{نعم} \\
C & تجهيز المعدات & 6 & 7 & 1 & لا \\
D & جمع البيانات & 10 & 10 & 0 & \color{criticalColor}\textbf{نعم} \\
E & التحليل الإحصائي & 20 & 20 & 0 & \color{criticalColor}\textbf{نعم} \\
F & النمذجة والمحاكاة & 20 & 21 & 1 & لا \\
G & التحقق من النتائج & 26 & 27 & 0 & \color{criticalColor}\textbf{نعم} \\
H & كتابة الورقة & 30 & 31 & 1 & لا \\
I & تقديم للنشر & 38 & 39 & 1 & لا \\
J & التقرير النهائي & 30 & 30 & 0 & \color{criticalColor}\textbf{نعم} \\
\bottomrule
\end{tabularx}
\end{table}

\textbf{ملاحظات:}
\begin{itemize}
    \item المهام الحرجة (السماحية = 0) يجب أن تُنجز في وقتها تماماً، أي تأخير فيها يؤخر المشروع كله.
    \item المهام غير الحرجة لها سماحية (\LR{slack}) يمكن استغلالها.
    \item يجب مراقبة المسار الحرج باستمرار أثناء تنفيذ المشروع.
\end{itemize}

\end{document}
